% s-061-p061-m003; original PDF 61, 62
\Aphorism{s-061-p061-m003}{60}
Wenn wir einen grossen Mann verehren, sehen wir eine Lichtgestalt, wo nur zwei oder drei Lichtflecke sind; unsere constructive Phantasie ergänzt das Nöthige. Wenn wir ein Weib lieben, stellen wir uns etwas Alabasterweisses und Marmorglattes vor, wo bei genauerem Hinschauen Dunst und Poren und behaarte Mäler sind. Diese Art Blindheit und Idealismus, dieses Wahrnehmen, das eigentlich ein Falschnehmen ist, fehlt uns im Verkehre mit uns selbst; daher ist der Egoismus leichter gesagt als gethan. Es ist schwer, sich so gut zu stilisiren, dass man sich selber wohlgefällt; es ist schwer, sich vor sich selber so in Scene zu setzen, dass man sich nicht unter die eigene Maske guckt; es ist schwer, vor Jemandem seine Hochachtung zu bewahren, mit dem man in der indiscretesten aller Beziehungen, der Identität, steht. Liebe deinen Nächsten wie dich selbst — aber man wird ihn immer mehr lieben; man wird immer die erste Person gegen die zweite und die zweite gegen die dritte geringschätzen, die blaue Ferne wird jederzeit mehr gelten als die nächste Nachbarschaft. Mit wachsender Vereinigung der Liebenden erlischt die Liebe: wie kann man lieben, was man \Sperr{sein} muss?
\par

